\section{Introduction}

Cosmological structure formation describes how primordial quantum fluctuations generated during inflation evolve under gravitational instability into the complex cosmic web observed today.
While linear perturbation theory provides an accurate description of density perturbations at early epochs and large spatial scales, late-time evolution on intermediate and non-linear scales requires specialized analytical approximations and numerical frameworks.

This report documents the methodology and numerical validation for the Numerical Cosmology II (EP2) project. The developed \texttt{Python} framework integrates six core cosmological pipelines:
    \begin{enumerate}
        \item \textbf{Background FLRW Dynamics \& Linear Growth}: High-precision numerical ODE integration of the growth factor $D(a)$ and logarithmic growth rate $f(a)$ across flat $\Lambda$CDM, Einstein--de Sitter (EdS), and open CDM cosmologies.
        
        \item \textbf{Gaussian Random Fields (GRF)}: Fourier-space generation of 3D linear density realizations enforcing Hermitian symmetry and exact mass conservation.
        
        \item \textbf{Standard Perturbation Theory (SPT)}: Real-space FFT evaluation of second-order density fields $\delta_2(\mathbf{x})$, equilateral bispectrum $B(k,k,k)$, and Top-Hat smoothed reduced skewness $S_3(R)$.
        
        \item \textbf{Lagrangian Perturbation Theory (1LPT)}: Particle displacement fields under the Zel'dovich approximation, Nearest Grid Point (NGP) mass assignment, and shot-noise-subtracted 1D power spectrum measurement.
        
        \item \textbf{Spherical Collapse Solver}: Non-linear integration across three algebraically equivalent formulations (Models A, B, and C) and bisection shooting determination of the linear collapse threshold $\delta_c(z)$.
        
        \item \textbf{Boltzmann Map Integration}: Transfer function reconstruction and cross-code precision benchmarking between CLASS and CAMB linear matter power spectra $P_m(k)$.
    \end{enumerate}
    
As the main reference, we follow the lecture notes from the Numerical Cosmology course \parencite{Tiago_Batalha_Lecture_5, Tiago_Batalha_Lecture_6, Tiago_Batalha_Lecture_7, Tiago_Batalha_Lecture_8, Tiago_Batalha_Lecture_9}. 